\chapter{Methods}

\section{Basics}

\subsection{GNN libraries}

\subsubsection{Pytorch Geometric}

Pytorch Geometric provides the most common GNN models, e.g. GCNs and GINs. These types are usually implemented in two versions.
Pytorch \cite{Paszke.03.12.2019}

\subsubsection{DGL}

The DGL Model Zoo encompasses many GNN architectures for molecular property prediction, including neural versions of fingerprints (currently: AttentiveFP Predictor, GAT Predictor, GATv2 Predictor, MGCN Predictor, MPNN Predictor, SchNet Predictor, Weave Predicot, GIN Predictor, GNN OGB Predictor, Neural Fingerprint Predictor, Path-Augmented Grpah Transformer Predictor).

\subsubsection{PyGHO}

Pytorch Geometric Higher Order (PyGHO) provides models and appropriate data structures for Higher-Order GNNs \cite{Wang.28.11.2023}.

It should be noted that, in contrast to Pytorch Geometric and DGL, this recently published library was not heavily used until now; thus it cannot be ascertained that the implementations are fully correct and actually represent the proposed models.

\section{Evaluation and analysis}

\subsection{Tanimoto similarity}

Tanimoto similarity is a commonly used metric for assessing the similarity between molecular compound.
The general formula is equivalent to Jaccard similarity and given by
\[
\text{Tanimoto / Jaccard Similarity} = \frac{|A \cap B|}{|A \cup B|}
\]


In a first step, the Morgan fingerprint of the molecules of interest is computed.
(It should be noted that the settings for the Morgan fingerprint can influence the results, especially a fingerprint size set much too small leading to collisions.) Afterwards, the similarity between the two fingerprint vectors is calculated.





In \cite{Duvenaud.30.09.2015}, the following generalization for continuous data is used for comparing neural and classical circular fingerprints:
\[
\text{distance}(x, y) = 1 - \frac{\sum \min(x_k, y_k)}{\sum \max(x_k, y_k)}
\]


\subsection{Metrics}



\subsubsection{Regression tasks}

The most commonly used loss functions are MAE and MSE loss, for the error of the final prediction either RMSE, 

\begin{equation}
\text{RMSE} = \sqrt{\frac{1}{n} \sum_{i=1}^{n} (y_i - \hat{y}_i)^2}
\end{equation}
and MAE, 
\begin{equation}
\text{MAE} = \frac{1}{n} \sum_{i=1}^{n} |y_i - \hat{y}_i|
\end{equation}
is used.
This could be of interest for the tasks since often outliers influence the results dramatically.


\subsubsection{Classification tasks}


For benchmarking of the molecular tasks, often AUROC is used. However, due to the imbalance of most datasets, its usage has also been rightly criticized (\cite{Robinson.2020}) as misleading (at least when it is used as performance benchmark without combination of other metrics).

In the following experiments for classification tasks primarily AUPRC is used.

Precision or positive predictive value (PPV) is defined as
\begin{equation}
\text{Precision} = \frac{\text{True Positives}}{\text{True Positives} + \text{False Positives}}
\end{equation}

\begin{equation}
\text{Recall} = \frac{\text{True Positives}}{\text{True Positives} + \text{False Negatives}}
\end{equation}

Alternatively, in \cite{Jiang.2021}, positive predictive (PPV), i.e. precision, in combination with negative predictive value (NPV),

\begin{equation}
\text{NPV} = \frac{\text{True Negatives}}{\text{True Negatives} + \text{False Negatives}}
\end{equation}
, is recommended.


Since PPV and NPV require a pre-defined threshold, it seems that AUPRC is most convenient for comparing the performance across different architectures.
